\documentclass{article}
\usepackage{mathtools} 
\usepackage{fontspec}
\usepackage[UTF8]{ctex}
\usepackage{amsthm}
\usepackage{mdframed}
\usepackage{xcolor}
\usepackage{amssymb}
\usepackage{amsmath}


% 定义新的带灰色背景的说明环境 zremark
\newmdtheoremenv[
  backgroundcolor=gray!10,
  % 边框与背景一致，边框线会消失
  linecolor=gray!10
]{zremark}{说明}

% 通用矩阵命令: \flexmatrix{矩阵名}{元素符号}{行数}{列数}
\newcommand{\flexmatrix}[4]{
  \[
  #1 = \begin{pmatrix}
    #2_{11}     & #2_{12}     & \cdots & #2_{1#4}   \\
    #2_{21}     & #2_{22}     & \cdots & #2_{2#4}   \\
    \vdots      & \vdots      & \ddots & \vdots     \\
    #2_{#31}    & #2_{#32}    & \cdots & #2_{#3#4}
  \end{pmatrix}
  \]
}

% 简化版命令（默认矩阵名为A，元素符号为a）: \quickmatrix{行数}{列数}
\newcommand{\quickmatrix}[2]{\flexmatrix{A}{a}{#1}{#2}}

\begin{document}
\title{注释 5.3}
\author{张志聪}
\maketitle
\begin{zremark}
  \begin{align*}
    \frac{dy}{dx} = \frac{\psi'(t)}{\varphi'(t)}
  \end{align*}
  的详细推导过程。
\end{zremark}

\textbf{证明：}

我们有
\begin{align*}
  t = \varphi^{-1}(x)
\end{align*}
按照定理5.7（反函数的导数），我们有
\begin{align*}
  \frac{dx}{dt} & = x'(t)                        \\
                & = \varphi'(t)                  \\
                & = \frac{1}{(\varphi^{-1}(x))'}
\end{align*}
即
\begin{align*}
  (\varphi^{-1}(x))' = \frac{1}{\varphi'(t)}
\end{align*}
$y$表示成关于$x$的函数，
\begin{align*}
  y & = \psi(t)               \\
    & = \psi(\varphi^{-1}(x))
\end{align*}
按照复合函数的求导公式，我们有
\begin{align*}
  \frac{dy}{dx}
   & = \psi'(\varphi^{-1}(x))(\varphi^{-1}(x))' \\
   & = \psi'(t)(\varphi^{-1}(x))'               \\
   & = \frac{\psi'(t)}{\varphi'(t)}
\end{align*}


\end{document}